% Resumo — Manual FATEC 4.1.10 (parágrafo único, espaçamento simples,
% palavras-chave separadas e finalizadas por ponto).
\begin{resumo}
O cenário econômico brasileiro recente, caracterizado por instabilidade
macroeconômica, inflação elevada e consequente aumento da taxa Selic, impõe
desafios significativos aos investidores. A volatilidade do mercado financeiro,
intensificada por fatores internos e externos, dificulta a tomada de decisões de
investimento, especialmente em ativos como ações. Diante desse contexto e da
crescente digitalização do setor financeiro, o trading algorítmico apresenta-se como
uma alternativa promissora. Este trabalho apresenta o desenvolvimento de um algoritmo de
negociação, escrito na linguagem MQL5 para a plataforma MetaTrader 5, para
análise de padrões de tendência e volatilidade no mercado de ações brasileiro,
tomando como base o ativo PETR4. O algoritmo integra análises técnicas
estatísticas, combinando o uso de Médias Móveis para identificação de tendências
e os modelos ARIMA para previsão de séries temporais de preços e GARCH para
modelagem da volatilidade condicional dos resíduos do ARIMA. A premissa central
deste trabalho é que essa abordagem integrada permite a geração de sinais de
negociação com maior eficácia (potencial de acerto e adequação ao risco) em
comparação com o uso isolado ou incompleto dessas técnicas. O estudo explorou
modelos com dados atualizados, fornecendo análises que possam embasar estratégias de
investimento adaptadas ao cenário de juros altos e investigando a robustez da
combinação de modelos ARIMA e GARCH frente às flutuações do mercado, contribuindo para o
conhecimento na área e oferecendo uma alternativa prática para investidores.

\vspace{\onelineskip}
\noindent\textbf{Palavras-chave:} Trading Algorítmico. Modelos Estatísticos. ARIMA.
GARCH. Mercado de Ações. Bolsa de Valores. Volatilidade. MetaTrader 5.
\end{resumo}
